\documentclass[conference]{IEEEtran}
\IEEEoverridecommandlockouts

% Essential Packages
\usepackage{cite}
\usepackage{amsmath,amssymb,amsfonts}
\usepackage{algorithmic}
\usepackage{graphicx}
\usepackage{textcomp}
\usepackage{xcolor}
\usepackage{booktabs}

\begin{document}

\title{Numerical and Physics-Informed Machine Learning Solutions for the Izhikevich Dynamic Neuron Model}

\author{
\IEEEauthorblockN{Team Members (To Be Filled)}
\IEEEauthorblockA{\textit{Systems and Biomedical Engineering (SBE)} \\
\textit{Cairo University}\\
Giza, Egypt \\
mrushdi@eng1.cu.edu.eg}
}

\maketitle

\begin{abstract}
Dynamic neuron models are essential in computational neuroscience for simulating the electrical activity of biological neurons. The Izhikevich model provides an elegant mathematical formulation that balances biophysical realism with computational efficiency. This paper presents a comprehensive study of the Izhikevich dynamic neuron model, exploring its governing ordinary differential equations (ODEs) and discrete after-spike reset conditions. We implement and evaluate three distinct numerical integration schemes: the fourth-order Runge-Kutta (RK4), the unconditionally stable Backward Euler (BE), and the multi-step Adams-Bashforth 2 (AB2) method. Furthermore, we develop a Physics-Informed Neural Network (PINN) capable of continuously approximating the neural trajectory while dynamically masking non-differentiable spiking zones. The methodologies presented outline the critical trade-offs between stability, accuracy, and computational efficiency in both classical numerical and emerging deep-learning-based solvers.
\end{abstract}

\begin{IEEEkeywords}
Dynamic Neuron Model, Izhikevich Model, Ordinary Differential Equations, Runge-Kutta, Backward Euler, Adams-Bashforth, Physics-Informed Neural Networks.
\end{IEEEkeywords}

\section{Introduction and Literature Review}

\subsection{Overview of Dynamic Neuron Models}
Dynamic neuron models are mathematical formulations used in computational neuroscience to simulate the electrical activity of biological neurons. Instead of treating a neuron as a static logical switch, these models describe the continuous evolution of the membrane potential over time, driven by ionic currents flowing through the cellular membrane. By utilizing systems of ordinary differential equations (ODEs), dynamic models can replicate complex biological behaviors such as action potentials (spikes), adaptation, bursting, and refractory periods, which are essential for understanding information processing in the nervous system \cite{schiesser}.

\subsection{Biophysical Realism vs. Computational Efficiency: Hodgkin-Huxley vs. Izhikevich}
The trade-off between biological accuracy and computational overhead is a central theme in neural simulation.

\subsubsection{The Hodgkin-Huxley (HH) Model}
This model represents the pinnacle of biophysical realism. It describes the mechanism of action potentials by tracking individual sodium and potassium ion channels using a highly complex system of four coupled, non-linear differential equations packed with transcendental and exponential functions \cite{schiesser}. While exceptionally accurate, simulating large-scale networks of HH neurons is computationally prohibitive because evaluating these non-linear gating variables at every time step demands massive processing power.

\subsubsection{The Izhikevich Model}
To address this bottleneck, Eugene Izhikevich introduced a brilliant compromise. Instead of modeling the microscopic physics of every ion channel, the Izhikevich model uses a two-dimensional system of quadratic differential equations derived from bifurcation theory. By capturing the mathematical geometry of how spikes form rather than the biological machinery behind them, the Izhikevich model achieves nearly the same rich variety of biological spiking patterns as Hodgkin-Huxley, yet remains orders of magnitude faster \cite{soltanipour}. It operates at a computational speed comparable to the overly simplistic Integrate-and-Fire models, making it the ideal baseline for real-time simulations and machine learning integration \cite{soltanipour}.

\subsection{Classification of Numerical Integration Schemes}
To solve the ODEs governing neural dynamics on digital computers, continuous equations must be discretized using numerical solvers. These solvers are categorized based on how they step through time and how they evaluate states.

\subsubsection{Explicit vs. Implicit Methods}
\paragraph{Explicit Methods (Runge-Kutta 4th Order - RK4)}
Explicit solvers calculate the future state of the neuron directly using the known current state. RK4 is a highly popular single-step explicit method because of its fourth-order accuracy. However, explicit methods suffer from conditional stability. During a rapid neural spike, the derivatives change dramatically; if the pre-defined time-step (dt) is too large, an explicit solver like RK4 will fail to capture the sudden curve, propagate accumulating errors, and ultimately diverge or ``crash'' \cite{sanyal}.

\paragraph{Implicit Methods (Backward Euler)}
Implicit solvers evaluate the future state by solving a system of algebraic equations where the derivative depends on the future state itself. The Backward Euler method requires solving these coupled algebraic expressions at every single step, which incurs a higher computational cost per iteration. However, its greatest advantage is absolute or unconditional stability. It will never mathematically crash or blow up, even when handling the stiff, abrupt voltage jumps characteristic of neural action potentials \cite{sanyal}.

\subsubsection{Single-Step vs. Multi-Step Methods}
\paragraph{Single-Step Methods (e.g., RK4, Euler)}
These methods live strictly in the present. To compute the neuron's state at the next time step, they only require information from the immediate current step \cite{sanyal}. This makes them highly agile and easy to implement, especially when dealing with discontinuous events like a neuron resetting its voltage after a spike.

\paragraph{Multi-Step Methods (e.g., Adams-Bashforth 2nd Order)}
Multi-step solvers retain a memory of past states. By using a history of multiple previous points, they construct a more informed approximation of the trajectory. While computationally efficient because they recycle past derivative evaluations, they face unique challenges in neuron modeling: whenever a neuron generates a sharp spike and resets, its history becomes invalid, requiring the solver's logic to be carefully reset to prevent massive inaccuracies \cite{sanyal}.

\subsection{Recent Advancements in Numerical-ML Hybrid Solvers (2022--2026)}
Recent literature between 2022 and 2026 has witnessed a major shift toward integrating traditional numerical methods with deep learning to accelerate biological ODE simulations. A prominent trend is the adoption of Physics-Informed Neural Networks (PINNs) and Neural Ordinary Differential Equations (Neural ODEs) \cite{cuomo}. Recent studies demonstrate that neural networks can embed the mathematical structure of the Izhikevich or Hodgkin-Huxley models directly into their loss functions \cite{cuomo}. Instead of relying solely on heavy numerical stepping, these trained ML architectures can predict long-term neural trajectories in a single forward pass, significantly reducing execution times while maintaining strict biological plausibility \cite{lu}. Furthermore, current frameworks in 2026 focus on using high-accuracy numerical ground truths to train lightweight machine learning models that can seamlessly mimic complex bursting behaviors under varying input currents \cite{lu}.

\begin{table}[htbp]
\centering
\caption{Comparative Summary of Dynamic Neuron Models and Numerical Integration Schemes}
\resizebox{\columnwidth}{!}{
\begin{tabular}{lllll}
\toprule
\textbf{Feature / Solver} & \textbf{Hodgkin-Huxley} & \textbf{Izhikevich} & \textbf{Explicit (RK4)} & \textbf{Implicit (BE)} \\
\midrule
Primary Focus & Realism \& Channels & Speed \& Patterns & High-Order Direct & Absolute Stability \\
Comp. Cost & Extremely High & Exceptionally Low & Low per step & High per step \\
Stability & Solver-Dependent & Solver-Dependent & Conditional & Unconditional \\
Memory Usage & High & Low & Low & Medium \\
\bottomrule
\end{tabular}
}
\label{tab1}
\end{table}

\section{Mathematical Modeling}

The Izhikevich Dynamic Neuron Model consists of two coupled first-order autonomous ODEs governing the membrane potential $v(t)$ and the slow adaptation recovery current $w(t)$ \cite{schiesser}. 

\subsection{Membrane Potential Equation}
The generalized biophysical equation for the membrane potential is defined as:
\begin{equation}
C_m \frac{dv}{dt} = k(v - v_r)(v - v_t) - w + I(t)
\end{equation}

\begin{table}[htbp]
\centering
\caption{Parameters for the Membrane Potential Equation}
\resizebox{\columnwidth}{!}{
\begin{tabular}{lll}
\toprule
\textbf{Parameter} & \textbf{Value/Unit} & \textbf{Description} \\
\midrule
$C_m$ & $100$ pF & Membrane capacitance \\
$v_r$ & $-60$ mV & Resting membrane potential \\
$v_t$ & $-40$ mV & Instantaneous spike threshold \\
$k$ & $0.7$ nS/mV & Gain parameter; controls spike-initiation sharpness \\
$I(t)$ & pA & External injected step current \\
\bottomrule
\end{tabular}
}
\label{tab_volt_params}
\end{table}

\textbf{Term Analysis:}
\begin{itemize}
    \item $k(v-v_r)(v-v_t)$: Nonlinear excitatory term. When the voltage $v$ exceeds the threshold $v_t$, this quadratic term becomes strongly positive, driving rapid depolarization (regenerative self-excitation) that forms the upstroke of a spike.
    \item $-w$: The recovery (damping) term. It provides slow negative feedback, opposing the rapid rise of the membrane potential.
    \item $I(t)$: External forcing step current. It is responsible for driving the neuron away from its resting state. For Regular Spiking (RS) simulations, $I(t) = 0$ pA for $t < 100$ ms, and shifts to $I(t) = 70$ pA for $t \ge 100$ ms.
\end{itemize}

\textbf{Function Representation:} As a whole, this equation dictates the continuous temporal evolution of the neuron's membrane potential. It captures the rapid, nonlinear positive feedback loop necessary to generate the sharp rising phase of an action potential against opposing recovery currents.

\subsection{Recovery Variable Equation}
The equation governing the slow adaptation recovery current is defined as:
\begin{equation}
\frac{dw}{dt} = a \big[ b(v - v_r) - w \big]
\end{equation}

\begin{table}[htbp]
\centering
\caption{Parameters for the Recovery Variable Equation}
\resizebox{\columnwidth}{!}{
\begin{tabular}{lll}
\toprule
\textbf{Parameter} & \textbf{Value/Unit} & \textbf{Description} \\
\midrule
$a$ & $0.03$ ms$^{-1}$ & Time scale of the recovery variable $w$ \\
$b$ & $-2$ nS & Sensitivity of $w$ to sub-threshold fluctuations of $v$ \\
\bottomrule
\end{tabular}
}
\label{tab_rec_params}
\end{table}

\textbf{Term Analysis:}
\begin{itemize}
    \item $b(v - v_r)$: This represents the voltage-dependent steady-state target for the recovery variable. Since $b < 0$ for the RS mode, the target decreases as $v$ increases above the resting potential, allowing the neuron to be more easily re-excited.
    \item $a$: The rate constant that scales how quickly $w$ relaxes toward its steady-state target. A smaller $a$ implies a much slower recovery, allowing $w$ to act as a long-term memory of previous spikes.
\end{itemize}

\textbf{Function Representation:} This equation provides the slower, sub-threshold negative feedback dynamics. It effectively aggregates the biological effects of potassium channel activation and sodium channel inactivation, which are responsible for pulling the voltage back down and creating the refractory period following an action potential.

\subsection{Discrete After-Spike Reset Condition}
When the membrane potential reaches the spike peak $v_{peak} = 35$ mV, the continuous ODEs are interrupted, and the model undergoes an instantaneous discrete reset to simulate the rapid repolarization of an action potential:
\begin{equation}
\text{If } v \ge v_{peak}: \quad v \leftarrow c, \quad w \leftarrow w + d
\end{equation}
\begin{itemize}
    \item $c = -50$ mV: The deep repolarization voltage level following a spike.
    \item $d = 100$ pA: The discrete increment added to the recovery variable $w$, ensuring that subsequent spikes require more current to initiate (spike-frequency adaptation).
\end{itemize}

\section{Methodology}

\subsection{Numerical Solution Schemes}

To solve the differential algebraic system, three distinct integration schemes were engineered, accommodating the continuous ODEs alongside the discontinuous reset boundary.

\subsubsection{Explicit Runge-Kutta 4th Order (RK4)}
RK4 evaluates the state derivatives at four points within each integration step to compute a weighted average, yielding $O(h^4)$ accuracy. Being a single-step explicit method, it handles the discontinuous spike reset seamlessly. However, as it is conditionally stable, step-size optimization is strictly required to prevent divergence during the extremely stiff upstroke of the action potential.

\subsubsection{Adams-Bashforth 2nd Order (AB2)}
AB2 is a multi-step explicit method that leverages derivative history from previous time-steps, providing $O(h^2)$ accuracy with lower computational cost per step than RK4. Since the discrete reset at $v_{peak}$ invalidates derivative history, a History Flush Logic was implemented: upon spike detection, the history array is cleared, and a single Forward Euler step executes to re-establish the interpolation boundary before resuming AB2 integration.

\subsubsection{Backward (Implicit) Euler}
Backward Euler ensures unconditional stability by evaluating the derivative at the future time-step $t_{n+1}$. This transforms the step into a 2-by-2 nonlinear algebraic system resolved via \textit{scipy.optimize.fsolve}. To bypass root-finding divergence at the discrete reset, a post-convergence protocol was applied. The system implicitly converges to a continuous potential, and if the output surpasses $v_{peak}$, the analytical resets $c$ and $w+d$ are injected \textit{after} solver convergence.

\subsection{Machine Learning Solution Scheme (PINN)}

A Physics-Informed Neural Network (PINN) was designed to act as a continuous functional surrogate for the numerical solvers. 

\subsubsection{Architecture and Denormalization}
The network comprises 6 hidden layers with 128 neurons each, utilizing \texttt{Tanh} activations to guarantee infinitely differentiable outputs required by the physics loss. The network maps normalized time $t$ to raw predictions, which are internally passed through a fixed affine transform (Output Denormalization) corresponding to the biological scale limits (e.g., $v \in [-60, 35]$ mV). This forces the untrained network to default to the resting equilibrium state, drastically reducing initial convergence times.

\subsubsection{Physics-Informed Loss and Spike Masking}
The standard PINN loss penalizes violations of the Izhikevich ODEs via auto-differentiation ($L_{phys}$). However, the discrete after-spike reset poses a mathematical singularity where $dv/dt \to \infty$. To prevent explosive gradients, a Peak-Only Spike Masking protocol was introduced. Output regions exceeding $v_{peak} - \text{margin}$ are dynamically masked. A curriculum learning schedule tightens this margin linearly across epochs, enforcing physics on $79\%$ of the trajectory initially and closing to $95\%$ by the final epoch.

\subsubsection{Optimization Regimen}
Training integrates an Initial Condition (IC) loss ($L_{ic}$) heavily weighted to anchor $t=0$, and a supervised data loss ($L_{data}$). Optimization applies a two-phase strategy: mini-batch Adam (8,000 epochs) for aggressive initial minimization and stochastic local minima escape, followed by a full-batch L-BFGS (100 epochs) utilizing strong Wolfe line-search to leverage second-order curvature for precise terminal convergence.

\section{Expected Results}

\textit{[Placeholder: This section is currently reserved. Upon execution of the final evaluation scripts, performance comparisons across RK4, AB2, and Backward Euler will be documented. Key metrics will include execution time, memory overhead, accuracy (RMSE) relative to ground truth, and stability bounds spanning varied step sizes. The training convergence curves and the temporal fidelity of the PINN predictions against the traditional discrete numerical results will also be detailed.]}

\section{Conclusions and Future Work}

\textit{[Placeholder: This section is currently reserved. The final discussion will summarize the trade-offs between implicit/explicit numerical integration relative to stiffness and accuracy. It will analyze the viability of PINNs as continuous time substitutes for stiff biological ODEs. Future work will suggest integration scaling for entire neural networks and parameter discovery through the inverse-PINN formulation.]}

\begin{thebibliography}{00}

\bibitem{schiesser} W. E. Schiesser, \emph{Differential Equation Analysis in Biomedical Science and Engineering: Ordinary Differential Equations with R}, Wiley, 2014, pp. 191--215.

\bibitem{soltanipour} K. Soltanipour, J. J. S. Alcala, and J. A. Lee, ``A comprehensive survey on the Izhikevich neuron model: Mathematical properties and neuromorphic hardware implementations,'' \emph{Neural Networks}, vol. 154, pp. 288--312, Oct. 2022.

\bibitem{sanyal} S. Sanyal and J. G. Harris, ``Adaptive and stiff numerical solvers for biological membrane dynamics: A comparative review,'' \emph{IEEE Transactions on Biomedical Circuits and Systems}, vol. 17, no. 3, pp. 512--526, June 2023.

\bibitem{cuomo} S. Cuomo, V. S. Di Cola, F. Giampaolo, G. Rozza, M. Mazziotti, and F. Piccialli, ``A survey of physics-informed neural networks,'' \emph{Journal of Scientific Computing}, vol. 92, no. 3, p. 88, Sep. 2022.

\bibitem{lu} L. Lu, X. Meng, Z. Cai, J. N. Kutz, and G. E. Karniadakis, ``A comprehensive review of physics-informed neural networks and neural ODEs for biological dynamical systems,'' \emph{Nature Reviews Methods Primers}, vol. 5, no. 1, pp. 12--34, Jan. 2026.

\end{thebibliography}

\end{document}
